\subsection{Database Implementation}

The database was implemented using PostgreSQL, with SQLAlchemy providing the Object-Relational Mapping (ORM) layer between the application and the database. Each table from the relational schema is represented by a corresponding ORM model defining primary keys, foreign keys and relationships.

Database schema changes are managed with Alembic migrations, allowing the schema to evolve in a controlled and reproducible manner without recreating the database. Frequently reused values, such as vehicle types, rental statuses, payment methods and fuel types, are stored in dedicated lookup tables to reduce redundancy and improve consistency.

Referential integrity is enforced through foreign key constraints, while additional business rules, including vehicle availability, rental validation and payment processing, are implemented in the backend before transactions are committed. Together with the deterministic seeding script, this provides a consistent and reproducible database environment for development and testing.
\newpage


\subsection{Backend Implementation}

The backend was implemented as a FastAPI application responsible for exposing the REST API, validating rental rules, communicating with the PostgreSQL database and keeping reservations, payments and invoices consistent. The code is divided into several main layers: \texttt{api} contains HTTP endpoints, \texttt{models} contains SQLAlchemy ORM classes, \texttt{schemas} contains Pydantic request and response models, \texttt{services} contains reusable business logic, and \texttt{core} contains infrastructure code such as configuration, security, database sessions, and time helpers.

The application entry point creates the FastAPI instance, configures CORS for the frontend and registers routers for authentication, customer operations, administrator operations and lookup data. Database access is handled through SQLAlchemy sessions injected into endpoints with FastAPI dependencies. Schema evolution and lookup-table seed data are managed with Alembic migrations.

\subsubsection{Database Communication}

Communication with the database is implemented with SQLAlchemy. The backend reads the database connection string from the \texttt{DATABASE\_URL} environment variable, creates a shared SQLAlchemy engine and then creates short-lived sessions for individual requests. Endpoints receive the session through the \texttt{get\_db} dependency, which ensures that the session is closed after the request is processed.

\begin{lstlisting}[style=pythonstyle,caption={Database engine and request-scoped session},label={lst:backend-db-short}]
DATABASE_URL = os.getenv("DATABASE_URL")

engine = create_engine(DATABASE_URL)
SessionLocal = sessionmaker(bind=engine, autoflush=False, autocommit=False)


def get_db():
    db = SessionLocal()
    try:
        yield db
    finally:
        db.close()
\end{lstlisting}

The same session is passed to service functions when more complex business logic is required. This allows a single operation, such as creating a reservation or processing a payment, to read data, validate relationships, insert records and commit changes as one transaction. For read and write operations the backend uses SQLAlchemy ORM models, while migrations and initial lookup data are managed separately with Alembic.

\subsubsection{Request and Response Schemas}

The API contract is separated from database models with Pydantic schemas. Request schemas define what data the frontend is allowed to send, while response schemas control which fields are returned to the client. This keeps SQLAlchemy models internal to the backend and provides an additional validation layer before data reaches business logic or database operations.

\subsubsection{Authentication and Authorization}

Authentication is based on JWT access tokens. Passwords are hashed with Argon2. Protected endpoints use dependencies to check the token payload. Customer routes require a customer account, while administrator routes require the \texttt{admin} role. This keeps access control close to the API layer and avoids duplicating role checks inside every endpoint.

\begin{lstlisting}[style=pythonstyle,caption={JWT payload validation and role guards},label={lst:backend-auth-short}]
def get_current_payload(credentials=Depends(bearer_scheme)) -> dict[str, Any]:
    try:
        return jwt.decode(credentials.credentials, SECRET_KEY, algorithms=[ALGORITHM])
    except InvalidTokenError:
        raise HTTPException(status_code=401, detail="Invalid or expired token")


def require_admin(payload=Depends(get_current_payload)) -> dict[str, Any]:
    if payload.get("role") != "admin":
        raise HTTPException(status_code=403, detail="Admin access required")
    return payload


def require_customer(payload=Depends(get_current_payload)) -> dict[str, Any]:
    if payload.get("account_type") != "customer":
        raise HTTPException(status_code=403, detail="Customer access required")
    return payload
\end{lstlisting}

\subsubsection{Customer Registration and Validation}

Customer registration creates both the customer account and the related address record. Before saving data, the backend normalizes the e-mail address, checks whether it is already registered, hashes the password, validates the customer's age and verifies the driver's license data. The same validation service is reused later during profile updates and rental creation, which keeps customer-related rules consistent across the application.

\subsubsection{Reservation Flow}

The most important customer operation is creating a rental reservation. The endpoint validates the requested date range, rejects rentals in the past, checks the customer's age and driver's license, locks the selected car row, verifies that the car is not in an excluded status and checks whether there are conflicting rentals for the same period.

\begin{lstlisting}[style=pythonstyle,caption={Creating a rental reservation},label={lst:backend-rent-short}]
@router.post("/rent", response_model=RentalResponse, status_code=status.HTTP_201_CREATED)
def rent_car(payload: CarRentalRequest, current_customer=Depends(get_current_customer), db=Depends(get_db)) -> Rental:
    now = utc_now()
    # ... validate dates, customer age, driver's license and locations ...

    car = _get_locked_car_or_404(db, payload.car_id)
    if db.get(CarStatus, car.car_status_id).name in EXCLUDED_CAR_STATUSES:
        raise HTTPException(status_code=status.HTTP_409_CONFLICT, detail="Car is not available for rental")

    start_datetime, planned_end_datetime = normalize_rental_day_range(...)
    availability = check_car_availability_for_period(db, payload.car_id, start_datetime, planned_end_datetime, now=now)
    if not availability.is_available:
        raise HTTPException(status_code=status.HTTP_409_CONFLICT, detail="Car is not available")

    rental = Rental(
        customer_id=current_customer.customer_id,
        car_id=payload.car_id,
        # ... locations, reserved status and date range ...
    )
    db.add(rental)
    db.commit()
    return rental
\end{lstlisting}

\subsubsection{Rental Lifecycle}

Rental state is controlled by a service module instead of being hard-coded in individual endpoints. A newly created unpaid rental receives the \texttt{reserved} status and blocks the car only for a short payment hold. Paid reservations block the period permanently, active rentals block the car immediately and expired unpaid reservations are cancelled automatically.

\begin{lstlisting}[style=pythonstyle,caption={Core availability rule for rentals},label={lst:backend-availability-short}]
ACTIVE_STATUS = "active"
RESERVATION_HOLD_MINUTES = 5


def check_car_availability_for_period(db: Session, car_id: UUID, start_date: datetime, planned_end_date: datetime) -> CarAvailabilityResult:
    now = utc_now()
    conflicts = db.execute(...).all()
    paid_ids = _paid_rental_ids(db, [rental.rental_id for rental, _ in conflicts])

    for rental, status_name in conflicts:
        if status_name == ACTIVE_STATUS or rental.rental_id in paid_ids:
            return CarAvailabilityResult(is_available=False)
        if is_reservation_hold_active(created_at=rental.created_at, now=now):
            return CarAvailabilityResult(is_available=False)
        # ... expired unpaid reservation is marked as cancelled ...

    return CarAvailabilityResult(is_available=True)
\end{lstlisting}

The same service also updates paid rentals according to dates: future rentals remain \texttt{reserved}, currently running rentals become \texttt{active}, finished rentals become \texttt{completed}. This makes the rental history reflect current business state.

\subsubsection{Payments and Invoices}

Payment finalizes a reservation. The backend verifies ownership of the rental, checks the payment method, rejects expired unpaid reservations, calculates the rental price from the number of days and the car's daily rate, applies an optional discount, creates an invoice and updates the rental lifecycle. Invoice PDF generation is implemented separately with ReportLab and is available as a downloadable \texttt{pdf} response.

\begin{lstlisting}[style=pythonstyle,caption={Payment and invoice creation},label={lst:backend-payment-short}]
@router.post("/payment", response_model=InvoiceResponse)
def pay_for_rental(payload: RentalPaymentRequest, current_customer=Depends(get_current_customer), db=Depends(get_db)) -> Invoice:
    rental = db.get(Rental, payload.rental_id)
    if rental is None or rental.customer_id != current_customer.customer_id:
        raise HTTPException(status_code=status.HTTP_404_NOT_FOUND, detail="Rental not found")

    now = utc_now()
    if not is_reservation_hold_active(created_at=rental.created_at, now=now):
        # ... mark rental as cancelled ...
        raise HTTPException(status_code=status.HTTP_409_CONFLICT, detail="Reservation expired")

    car = db.get(Car, rental.car_id)
    rental_days = count_rental_days(rental.start_date, rental.planned_end_date)
    base_amount = (car.daily_rate * rental_days).quantize(Decimal("0.01"))
    discount_amount = Decimal("0.00")
    # ... validate optional discount code and calculate discount_amount ...

    invoice = Invoice(
        rental_id=rental.rental_id,
        total_amount=base_amount - discount_amount,
        paid_at=now,
        # ... payment method, invoice number and statuses ...
    )

    db.add(invoice)
    db.commit()
    apply_paid_rental_lifecycle_status(db, rental, now=now)
    db.commit()
    return invoice
\end{lstlisting}

\subsubsection{Consistency}

The backend uses multiple consistency mechanisms: Pydantic schemas validate request payloads, SQLAlchemy transactions group related database changes, service functions enforce domain rules and database constraints catch duplicates or invalid relations. Invalid operations are translated into appropriate HTTP responses, for example \texttt{400 Bad Request} for validation errors, \texttt{404 Not Found} for missing records, \texttt{403 Forbidden} for unauthorized access, and \texttt{409 Conflict} for unavailable cars or expired reservations.


\subsection{Frontend Implementation}

The frontend of \projectTitle{} is a high-performance Single Page Application (SPA) built with \textbf{React 18} and \textbf{Vite}. The implementation emphasizes a modular component architecture, smooth visual feedback, and a clear separation between administrative and customer interfaces.




\subsubsection{Routing and Access Control}
The application utilizes \texttt{react-router-dom} for navigation. 
Access to specific parts of the system is managed through higher-order components called \textbf{Route Guards}:
\begin{itemize}
    \item \textbf{AppRouteGuard:} Restricts access to authenticated users based on their role (\texttt{customer} or \texttt{admin}).
    \item \textbf{GuestRoute:} Redirects authenticated users away from login and registration pages to their respective dashboards.
\end{itemize}


\subsubsection{Role-Based Functional Views}
The frontend differentiates between consumer-facing and management-facing logic:
\begin{itemize}
    \item \textbf{Customer Experience:} Focused on the \texttt{CarsPage}, which features a responsive grid of "Car Cards." Each card is interactive, allowing users to open a detailed modal or jump directly into the multi-step \texttt{UserRentPaymentPage} using React Router's \texttt{state} to pass car data.
    \item \textbf{Administrative Oversight:} Modules like \texttt{AdminUsersPage} provide CRUD (Create, Read, Update, Delete) capabilities. Administrative actions, such as user deletion, are protected by \texttt{ConfirmActionModal} to prevent accidental data loss.
\end{itemize}

The routing structure is divided into three primary modules:
\begin{enumerate}
    \item \textbf{Public:} Home, Login, and Signup pages.
    \item \textbf{Customer:} Car browsing (\texttt{/cars}), personal profile management, rental history, and the multi-step payment flow.
    \item \textbf{Admin:} Dashboard for fleet management, customer oversight, reservation tracking, and discount code configuration.
\end{enumerate}

\subsubsection{Core Layout and Component Architecture}
The application utilizes a "Shell" pattern to maintain a consistent user experience across different views. 
\begin{itemize}
    \item \textbf{PageShell:} A top-level wrapper that manages the primary navigation (\texttt{site-top-panel}) and global decorative elements. It uses role-based logic to conditionally render navigation links for \textit{Customers} (e.g., Profile, History, Cars) or \textit{Admins} (e.g., Users, Reservations, Coupons).
    \item \textbf{PageSection:} A structural component that standardizes page headers, providing consistent "eyebrows," titles, and action areas for all system modules.
    \item \textbf{Modal Orchestration:} The system uses specialized modals like \texttt{CarModal} and \texttt{EntityDetailsModal}. A notable implementation detail is the "Safe Value Unwrapping" logic, which ensures that complex data objects (such as nested JSON for car brands or fuel types) are correctly rendered as strings without causing React rendering errors.
\end{itemize}



\subsubsection{State Management and Session Persistence}
User sessions are managed using React's \texttt{useState} and \texttt{useEffect} hooks. 
Authentication tokens are persisted in local storage. 
The application monitors session expiration and includes a global event listener for \texttt{metrocars:unauthorized} 
events to automatically log users out if their token becomes invalid.

\subsubsection{Data Persistence and API Integration}
Communication with the FastAPI backend is centralized through a \texttt{requestJson} utility. This helper manages:
\begin{itemize}
    \item \textbf{JWT Injection:} Automatically attaches the Bearer token from the session state to every outgoing request.
    \item \textbf{Error Handling:} Provides fallback messaging and captures backend exceptions to update the UI's \texttt{error} state variables.
    \item \textbf{Loading States:} Every data-driven view (like \texttt{CarsPage} or \texttt{AdminUsersPage}) manages a \texttt{loading} boolean to show visual feedback during network latency.
\end{itemize}

\subsubsection{UI/UX Design and Visual Identity}
The interface employs a distinctive "Metro" theme characterized by a high-contrast palette (\texttt{\#2C2B2C} and \texttt{\#FFE75C}). Key visual features include:
\begin{itemize}
    \item \textbf{Dynamic Backdrop:} A \texttt{BoardBackdrop} component generates a procedurally placed grid of car SVGs and company logos that scales with the window size.
    \item \textbf{Motion Strips:} CSS-animated "motion strips" on the sides of the car browsing page simulate a road environment, enhancing the thematic consistency of the rental experience.
    \item \textbf{Interactive Components:} Custom-styled "button-pills" and "car-cards" provide a tactile feel with hover transformations and depth-based box shadows.
\end{itemize}


\subsubsection{Responsive Design and Theming}
The system uses a custom CSS architecture defined in \texttt{index.css}. 
It leverages CSS Variables (e.g., \texttt{--panel-accent}) and \texttt{backdrop-filter: blur()} to create a modern, layered glassmorphism effect. 
The application uses a mobile-first approach with CSS Grid and Flexbox. Using \texttt{clamp()} functions and media queries, 
the layout adapts from a multi-column desktop view (including the decorative motion strips) to a streamlined single-column view for mobile devices, 
ensuring accessibility across all platforms.


\subsubsection{Dynamic Visuals and Parallax Effects}
A key aesthetic feature of \projectTitle{} is the \texttt{CarMotionStrip}. Unlike static images, this component implements a scroll-driven parallax effect:
\begin{itemize}
    \item It utilizes \textbf{requestAnimationFrame} for high-performance 60fps animations.
    \item The system calculates the scroll position with a \texttt{0.35} multiplier to create a sense of depth as cars move vertically in the background.
    \item \textbf{GPU Acceleration:} The use of \texttt{translate3d(0, y, 0)} ensures that the animation is offloaded to the browser's hardware acceleration, preventing layout thrashing during heavy scrolling.
\end{itemize}

\begin{lstlisting}[style=pythonstyle, caption=High-Performance Scroll Animation Logic]
function updatePosition() {
  const loopHeight = track.scrollHeight / 2;
  // Calculate offset based on scroll position for parallax effect
  const travel = (window.scrollY * 0.35 * direction) % loopHeight;
  const offset = travel < 0 ? travel + loopHeight : travel;
  // Apply hardware-accelerated transform
  track.style.transform = `translate3d(0, ${-offset}px, 0)`;
}
\end{lstlisting}


\subsubsection{Interactive Features and Easter Eggs}
To enhance user engagement, the application includes hidden interactive elements. A specialized \textbf{MetroŻubr Easter Egg} is implemented via an invisible trigger on the interface's right margin, which opens a modal featuring unique brand-specific content. This demonstrates the use of React Portals and conditional rendering for non-intrusive UI additions.

\begin{lstlisting}[style=pythonstyle, caption=Snippet of App Routing and Guards]
<Routes>
  <Route path="/" element={<HomePage session={session} />} />
  <Route path="/admin/cars" element={
    <AppRouteGuard session={session} allowedRoles={["admin"]}>
      <AdminCarsPage />
    </AppRouteGuard>
  } />
  <Route path="/user/rent/payment" element={
    <AppRouteGuard session={session} allowedRoles={["customer"]}>
      <UserRentPaymentPage />
    </AppRouteGuard>
  } />
</Routes>
\end{lstlisting}

\section{Test Data}

To simplify development and provide a reproducible testing environment, the project includes a deterministic database seeding script (\texttt{seed.py}). The script automatically populates the database with representative records covering all major entities and their relationships.

The generated dataset consists of:

\begin{itemize}
    \item 12 rental locations distributed across cities in the Silesian region,
    \item 50 vehicles representing different categories, fuel types and transmission types,
    \item 100 customer accounts with complete personal information,
    \item 150 rentals covering completed, active and future reservations,
    \item invoices generated for rentals together with payment information,
    \item promotional discount codes,
    \item all lookup tables required by the application.
\end{itemize}

The lookup tables are initialized first and include predefined values for vehicle status, vehicle type, fuel type, transmission type, rental status, invoice status, payment status and payment methods.

Vehicles are generated using predefined model templates that specify attributes such as brand, model, fuel type, transmission, seating capacity and daily rental rate. Each vehicle is assigned a unique VIN number, registration plate, mileage and rental location.

Customer records contain realistic personal information including addresses, e-mail addresses, phone numbers and driving licence information. Passwords are automatically hashed using the Argon2 algorithm before being stored in the database.

Rental records are generated to represent different stages of the rental lifecycle. Completed rentals include invoices and payment information, active rentals occupy vehicles currently marked as rented, while future reservations simulate upcoming bookings. Rental prices are calculated automatically based on the vehicle's daily rate and rental duration. Selected rentals additionally receive promotional discounts.

Because the seed script is deterministic, every developer obtains exactly the same dataset after initialization, making application testing, demonstrations and debugging fully reproducible.

\subsection{Backend Implementation}

The backend was implemented as a FastAPI application responsible for exposing the REST API, validating rental rules, communicating with the PostgreSQL database and keeping reservations, payments and invoices consistent. The code is divided into several main layers: \texttt{api} contains HTTP endpoints, \texttt{models} contains SQLAlchemy ORM classes, \texttt{schemas} contains Pydantic request and response models, \texttt{services} contains reusable business logic, and \texttt{core} contains infrastructure code such as configuration, security, database sessions, and time helpers.

The application entry point creates the FastAPI instance, configures CORS for the frontend and registers routers for authentication, customer operations, administrator operations and lookup data. Database access is handled through SQLAlchemy sessions injected into endpoints with FastAPI dependencies. Schema evolution and lookup-table seed data are managed with Alembic migrations.

\subsubsection{Database Communication}

Communication with the database is implemented with SQLAlchemy. The backend reads the database connection string from the \texttt{DATABASE\_URL} environment variable, creates a shared SQLAlchemy engine and then creates short-lived sessions for individual requests. Endpoints receive the session through the \texttt{get\_db} dependency, which ensures that the session is closed after the request is processed.

\begin{lstlisting}[style=pythonstyle,caption={Database engine and request-scoped session},label={lst:backend-db-short}]
DATABASE_URL = os.getenv("DATABASE_URL")

engine = create_engine(DATABASE_URL)
SessionLocal = sessionmaker(bind=engine, autoflush=False, autocommit=False)


def get_db():
    db = SessionLocal()
    try:
        yield db
    finally:
        db.close()
\end{lstlisting}

The same session is passed to service functions when more complex business logic is required. This allows a single operation, such as creating a reservation or processing a payment, to read data, validate relationships, insert records and commit changes as one transaction. For read and write operations the backend uses SQLAlchemy ORM models, while migrations and initial lookup data are managed separately with Alembic.

\subsubsection{Request and Response Schemas}

The API contract is separated from database models with Pydantic schemas. Request schemas define what data the frontend is allowed to send, while response schemas control which fields are returned to the client. This keeps SQLAlchemy models internal to the backend and provides an additional validation layer before data reaches business logic or database operations.

\subsubsection{Authentication and Authorization}

Authentication is based on JWT access tokens. Passwords are hashed with Argon2. Protected endpoints use dependencies to check the token payload. Customer routes require a customer account, while administrator routes require the \texttt{admin} role. This keeps access control close to the API layer and avoids duplicating role checks inside every endpoint.

\begin{lstlisting}[style=pythonstyle,caption={JWT payload validation and role guards},label={lst:backend-auth-short}]
def get_current_payload(credentials=Depends(bearer_scheme)) -> dict[str, Any]:
    try:
        return jwt.decode(credentials.credentials, SECRET_KEY, algorithms=[ALGORITHM])
    except InvalidTokenError:
        raise HTTPException(status_code=401, detail="Invalid or expired token")


def require_admin(payload=Depends(get_current_payload)) -> dict[str, Any]:
    if payload.get("role") != "admin":
        raise HTTPException(status_code=403, detail="Admin access required")
    return payload


def require_customer(payload=Depends(get_current_payload)) -> dict[str, Any]:
    if payload.get("account_type") != "customer":
        raise HTTPException(status_code=403, detail="Customer access required")
    return payload
\end{lstlisting}

\subsubsection{Customer Registration and Validation}

Customer registration creates both the customer account and the related address record. Before saving data, the backend normalizes the e-mail address, checks whether it is already registered, hashes the password, validates the customer's age and verifies the driver's license data. The same validation service is reused later during profile updates and rental creation, which keeps customer-related rules consistent across the application.

\subsubsection{Reservation Flow}

The most important customer operation is creating a rental reservation. The endpoint validates the requested date range, rejects rentals in the past, checks the customer's age and driver's license, locks the selected car row, verifies that the car is not in an excluded status and checks whether there are conflicting rentals for the same period.

\begin{lstlisting}[style=pythonstyle,caption={Creating a rental reservation},label={lst:backend-rent-short}]
@router.post("/rent", response_model=RentalResponse, status_code=status.HTTP_201_CREATED)
def rent_car(payload: CarRentalRequest, current_customer=Depends(get_current_customer), db=Depends(get_db)) -> Rental:
    now = utc_now()
    # ... validate dates, customer age, driver's license and locations ...

    car = _get_locked_car_or_404(db, payload.car_id)
    if db.get(CarStatus, car.car_status_id).name in EXCLUDED_CAR_STATUSES:
        raise HTTPException(status_code=status.HTTP_409_CONFLICT, detail="Car is not available for rental")

    start_datetime, planned_end_datetime = normalize_rental_day_range(...)
    availability = check_car_availability_for_period(db, payload.car_id, start_datetime, planned_end_datetime, now=now)
    if not availability.is_available:
        raise HTTPException(status_code=status.HTTP_409_CONFLICT, detail="Car is not available")

    rental = Rental(
        customer_id=current_customer.customer_id,
        car_id=payload.car_id,
        # ... locations, reserved status and date range ...
    )
    db.add(rental)
    db.commit()
    return rental
\end{lstlisting}

\subsubsection{Rental Lifecycle}

Rental state is controlled by a service module instead of being hard-coded in individual endpoints. A newly created unpaid rental receives the \texttt{reserved} status and blocks the car only for a short payment hold. Paid reservations block the period permanently, active rentals block the car immediately and expired unpaid reservations are cancelled automatically.

\begin{lstlisting}[style=pythonstyle,caption={Core availability rule for rentals},label={lst:backend-availability-short}]
ACTIVE_STATUS = "active"
RESERVATION_HOLD_MINUTES = 5


def check_car_availability_for_period(db: Session, car_id: UUID, start_date: datetime, planned_end_date: datetime) -> CarAvailabilityResult:
    now = utc_now()
    conflicts = db.execute(...).all()
    paid_ids = _paid_rental_ids(db, [rental.rental_id for rental, _ in conflicts])

    for rental, status_name in conflicts:
        if status_name == ACTIVE_STATUS or rental.rental_id in paid_ids:
            return CarAvailabilityResult(is_available=False)
        if is_reservation_hold_active(created_at=rental.created_at, now=now):
            return CarAvailabilityResult(is_available=False)
        # ... expired unpaid reservation is marked as cancelled ...

    return CarAvailabilityResult(is_available=True)
\end{lstlisting}

The same service also updates paid rentals according to dates: future rentals remain \texttt{reserved}, currently running rentals become \texttt{active}, finished rentals become \texttt{completed}. This makes the rental history reflect current business state.

\subsubsection{Payments and Invoices}

Payment finalizes a reservation. The backend verifies ownership of the rental, checks the payment method, rejects expired unpaid reservations, calculates the rental price from the number of days and the car's daily rate, applies an optional discount, creates an invoice and updates the rental lifecycle. Invoice PDF generation is implemented separately with ReportLab and is available as a downloadable \texttt{pdf} response.

\begin{lstlisting}[style=pythonstyle,caption={Payment and invoice creation},label={lst:backend-payment-short}]
@router.post("/payment", response_model=InvoiceResponse)
def pay_for_rental(payload: RentalPaymentRequest, current_customer=Depends(get_current_customer), db=Depends(get_db)) -> Invoice:
    rental = db.get(Rental, payload.rental_id)
    if rental is None or rental.customer_id != current_customer.customer_id:
        raise HTTPException(status_code=status.HTTP_404_NOT_FOUND, detail="Rental not found")

    now = utc_now()
    if not is_reservation_hold_active(created_at=rental.created_at, now=now):
        # ... mark rental as cancelled ...
        raise HTTPException(status_code=status.HTTP_409_CONFLICT, detail="Reservation expired")

    car = db.get(Car, rental.car_id)
    rental_days = count_rental_days(rental.start_date, rental.planned_end_date)
    base_amount = (car.daily_rate * rental_days).quantize(Decimal("0.01"))
    discount_amount = Decimal("0.00")
    # ... validate optional discount code and calculate discount_amount ...

    invoice = Invoice(
        rental_id=rental.rental_id,
        total_amount=base_amount - discount_amount,
        paid_at=now,
        # ... payment method, invoice number and statuses ...
    )

    db.add(invoice)
    db.commit()
    apply_paid_rental_lifecycle_status(db, rental, now=now)
    db.commit()
    return invoice
\end{lstlisting}

\subsubsection{Consistency}

The backend uses multiple consistency mechanisms: Pydantic schemas validate request payloads, SQLAlchemy transactions group related database changes, service functions enforce domain rules and database constraints catch duplicates or invalid relations. Invalid operations are translated into appropriate HTTP responses, for example \texttt{400 Bad Request} for validation errors, \texttt{404 Not Found} for missing records, \texttt{403 Forbidden} for unauthorized access, and \texttt{409 Conflict} for unavailable cars or expired reservations.


\subsection{Frontend Implementation}

The frontend of \projectTitle{} is a high-performance Single Page Application (SPA) built with \textbf{React 18} and \textbf{Vite}. The implementation emphasizes a modular component architecture, smooth visual feedback, and a clear separation between administrative and customer interfaces.




\subsubsection{Routing and Access Control}
The application utilizes \texttt{react-router-dom} for navigation. 
Access to specific parts of the system is managed through higher-order components called \textbf{Route Guards}:
\begin{itemize}
    \item \textbf{AppRouteGuard:} Restricts access to authenticated users based on their role (\texttt{customer} or \texttt{admin}).
    \item \textbf{GuestRoute:} Redirects authenticated users away from login and registration pages to their respective dashboards.
\end{itemize}


\subsubsection{Role-Based Functional Views}
The frontend differentiates between consumer-facing and management-facing logic:
\begin{itemize}
    \item \textbf{Customer Experience:} Focused on the \texttt{CarsPage}, which features a responsive grid of "Car Cards." Each card is interactive, allowing users to open a detailed modal or jump directly into the multi-step \texttt{UserRentPaymentPage} using React Router's \texttt{state} to pass car data.
    \item \textbf{Administrative Oversight:} Modules like \texttt{AdminUsersPage} provide CRUD (Create, Read, Update, Delete) capabilities. Administrative actions, such as user deletion, are protected by \texttt{ConfirmActionModal} to prevent accidental data loss.
\end{itemize}

The routing structure is divided into three primary modules:
\begin{enumerate}
    \item \textbf{Public:} Home, Login, and Signup pages.
    \item \textbf{Customer:} Car browsing (\texttt{/cars}), personal profile management, rental history, and the multi-step payment flow.
    \item \textbf{Admin:} Dashboard for fleet management, customer oversight, reservation tracking, and discount code configuration.
\end{enumerate}

\subsubsection{Core Layout and Component Architecture}
The application utilizes a "Shell" pattern to maintain a consistent user experience across different views. 
\begin{itemize}
    \item \textbf{PageShell:} A top-level wrapper that manages the primary navigation (\texttt{site-top-panel}) and global decorative elements. It uses role-based logic to conditionally render navigation links for \textit{Customers} (e.g., Profile, History, Cars) or \textit{Admins} (e.g., Users, Reservations, Coupons).
    \item \textbf{PageSection:} A structural component that standardizes page headers, providing consistent "eyebrows," titles, and action areas for all system modules.
    \item \textbf{Modal Orchestration:} The system uses specialized modals like \texttt{CarModal} and \texttt{EntityDetailsModal}. A notable implementation detail is the "Safe Value Unwrapping" logic, which ensures that complex data objects (such as nested JSON for car brands or fuel types) are correctly rendered as strings without causing React rendering errors.
\end{itemize}



\subsubsection{State Management and Session Persistence}
User sessions are managed using React's \texttt{useState} and \texttt{useEffect} hooks. 
Authentication tokens are persisted in local storage. 
The application monitors session expiration and includes a global event listener for \texttt{metrocars:unauthorized} 
events to automatically log users out if their token becomes invalid.

\subsubsection{Data Persistence and API Integration}
Communication with the FastAPI backend is centralized through a \texttt{requestJson} utility. This helper manages:
\begin{itemize}
    \item \textbf{JWT Injection:} Automatically attaches the Bearer token from the session state to every outgoing request.
    \item \textbf{Error Handling:} Provides fallback messaging and captures backend exceptions to update the UI's \texttt{error} state variables.
    \item \textbf{Loading States:} Every data-driven view (like \texttt{CarsPage} or \texttt{AdminUsersPage}) manages a \texttt{loading} boolean to show visual feedback during network latency.
\end{itemize}

\subsubsection{UI/UX Design and Visual Identity}
The interface employs a distinctive "Metro" theme characterized by a high-contrast palette (\texttt{\#2C2B2C} and \texttt{\#FFE75C}). Key visual features include:
\begin{itemize}
    \item \textbf{Dynamic Backdrop:} A \texttt{BoardBackdrop} component generates a procedurally placed grid of car SVGs and company logos that scales with the window size.
    \item \textbf{Motion Strips:} CSS-animated "motion strips" on the sides of the car browsing page simulate a road environment, enhancing the thematic consistency of the rental experience.
    \item \textbf{Interactive Components:} Custom-styled "button-pills" and "car-cards" provide a tactile feel with hover transformations and depth-based box shadows.
\end{itemize}


\subsubsection{Responsive Design and Theming}
The system uses a custom CSS architecture defined in \texttt{index.css}. 
It leverages CSS Variables (e.g., \texttt{--panel-accent}) and \texttt{backdrop-filter: blur()} to create a modern, layered glassmorphism effect. 
The application uses a mobile-first approach with CSS Grid and Flexbox. Using \texttt{clamp()} functions and media queries, 
the layout adapts from a multi-column desktop view (including the decorative motion strips) to a streamlined single-column view for mobile devices, 
ensuring accessibility across all platforms.


\subsubsection{Dynamic Visuals and Parallax Effects}
A key aesthetic feature of \projectTitle{} is the \texttt{CarMotionStrip}. Unlike static images, this component implements a scroll-driven parallax effect:
\begin{itemize}
    \item It utilizes \textbf{requestAnimationFrame} for high-performance 60fps animations.
    \item The system calculates the scroll position with a \texttt{0.35} multiplier to create a sense of depth as cars move vertically in the background.
    \item \textbf{GPU Acceleration:} The use of \texttt{translate3d(0, y, 0)} ensures that the animation is offloaded to the browser's hardware acceleration, preventing layout thrashing during heavy scrolling.
\end{itemize}

\begin{lstlisting}[style=pythonstyle, caption=High-Performance Scroll Animation Logic]
function updatePosition() {
  const loopHeight = track.scrollHeight / 2;
  // Calculate offset based on scroll position for parallax effect
  const travel = (window.scrollY * 0.35 * direction) % loopHeight;
  const offset = travel < 0 ? travel + loopHeight : travel;
  // Apply hardware-accelerated transform
  track.style.transform = `translate3d(0, ${-offset}px, 0)`;
}
\end{lstlisting}


\subsubsection{Interactive Features and Easter Eggs}
To enhance user engagement, the application includes hidden interactive elements. A specialized \textbf{MetroŻubr Easter Egg} is implemented via an invisible trigger on the interface's right margin, which opens a modal featuring unique brand-specific content. This demonstrates the use of React Portals and conditional rendering for non-intrusive UI additions.

\begin{lstlisting}[style=pythonstyle, caption=Snippet of App Routing and Guards]
<Routes>
  <Route path="/" element={<HomePage session={session} />} />
  <Route path="/admin/cars" element={
    <AppRouteGuard session={session} allowedRoles={["admin"]}>
      <AdminCarsPage />
    </AppRouteGuard>
  } />
  <Route path="/user/rent/payment" element={
    <AppRouteGuard session={session} allowedRoles={["customer"]}>
      <UserRentPaymentPage />
    </AppRouteGuard>
  } />
</Routes>
\end{lstlisting}

\section{Test Data}

To simplify development and provide a reproducible testing environment, the project includes a deterministic database seeding script (\texttt{seed.py}). The script automatically populates the database with representative records covering all major entities and their relationships.

The generated dataset consists of:

\begin{itemize}
    \item 12 rental locations distributed across cities in the Silesian region,
    \item 50 vehicles representing different categories, fuel types and transmission types,
    \item 100 customer accounts with complete personal information,
    \item 150 rentals covering completed, active and future reservations,
    \item invoices generated for rentals together with payment information,
    \item promotional discount codes,
    \item all lookup tables required by the application.
\end{itemize}

The lookup tables are initialized first and include predefined values for vehicle status, vehicle type, fuel type, transmission type, rental status, invoice status, payment status and payment methods.

Vehicles are generated using predefined model templates that specify attributes such as brand, model, fuel type, transmission, seating capacity and daily rental rate. Each vehicle is assigned a unique VIN number, registration plate, mileage and rental location.

Customer records contain realistic personal information including addresses, e-mail addresses, phone numbers and driving licence information. Passwords are automatically hashed using the Argon2 algorithm before being stored in the database.

Rental records are generated to represent different stages of the rental lifecycle. Completed rentals include invoices and payment information, active rentals occupy vehicles currently marked as rented, while future reservations simulate upcoming bookings. Rental prices are calculated automatically based on the vehicle's daily rate and rental duration. Selected rentals additionally receive promotional discounts.

Because the seed script is deterministic, every developer obtains exactly the same dataset after initialization, making application testing, demonstrations and debugging fully reproducible.
